% Step 4 synthesis: Updated compactness analysis with corrections
% S0910A03 Agent2 (Constructor) — incorporates TTCS results
%
% CORRECTIONS:
% 1. Strong maximum principle argument for |Φ|² > 0 is incorrect as stated.
%    At a zero, the Weitzenböck equation gives Δu ≤ 0 and Δu = -2|∇Φ|²,
%    forcing |∇Φ| = 0, but unique continuation for D_A Φ = 0 does not require
%    all second derivatives to vanish. The spinor CAN have zeros.
%    HOWEVER: in the exact equality limit (Prop 3.8), |Φ|² = const > 0 is
%    proven directly from the Hölder equality chain. So the Armstrong–LeBrun
%    contradiction holds if we route through Prop 3.8 rather than the max principle.
%
% 2. Non-collapsing does NOT follow from |Ric| ≤ Λ + Vol ≥ V₀ alone.
%    Counterexample: Σ_g × Σ_{g'} with degenerating hyperbolic metric on one factor
%    has exact Ric = -g, fixed Vol, but local collapse near the short geodesic.
%    BUT: this product has |W+| = 2/√6 (it's Kähler), which is excluded by
%    the Weyl profile.

\section{Corrected compactness chain}

\subsection{The argument (conditional on non-collapsing)}

Suppose \(\delta_i\to0\) on a fixed critical manifold \(M\) (Ishida--LeBrun,
\(D=A/3\)).  The chain is:

\begin{enumerate}
\item \(|\operatorname{Ric}_{g_i}|\le1+\delta_i\le2\) (uniform).
\item \(\operatorname{Vol}(M,g_i)\to32\pi^2A/9>0\) (from CGB + calibrated
  Weyl profile + SW near-equality).
  \emph{This uses the specific topology and monopole structure, not CGB alone.}
\item \textbf{[Gap]} Non-collapsing: \(\operatorname{Vol}(B(p,1))\ge v>0\)
  for all \(p\) and \(i\).
\item If (3) holds: Anderson's compactness gives \(C^{1,\alpha}\)
  subsequential convergence \(\varphi_i^*g_i\to g_\infty\)
  (modulo diffeomorphisms).  Need to also control curvature concentration
  (ALE bubbles); the Weyl profile's bubble restriction limits these to
  hyperk\"ahler models.  If no bubbles form, convergence is on all of \(M\).
\item \(\operatorname{Ric}_\infty=-g_\infty\) (distributional, from
  \(L^\infty\) convergence of Ricci).  DeTurck--Kazdan elliptic regularity
  in harmonic coordinates: \(g_\infty\) is smooth.
\item \(K(g_\infty)=(1/72\pi^2)\int(s_\infty-\sqrt6|W_{+,\infty}|)^2=A\)
  (from convergence of the integral identity and \(K_i\to A\)).
  By Proposition~3.8 (LeBrun 2006): equality forces
  \(|\Phi_\infty|=\text{const}\ne0\) directly (via H\"older equality),
  hence a smooth closed self-dual 2-form of constant nonzero norm.
  \emph{This does not use the strong maximum principle; the spinor
  non-vanishing follows from the equality characterization.}
\item Corollary~2.5 (Armstrong): Einstein + saturated almost-K\"ahler
  \(\Rightarrow\) K\"ahler--Einstein.  But then
  \(|W_+|^2=s^2/24=2/3\ne1/24\) (profile value).  Contradiction.
  Alternatively: \(g_\infty\) is an Einstein metric on \(M\), but
  Theorem~7 excludes all Einstein metrics on \(M\).  Double contradiction.
\end{enumerate}

\subsection{Correction: why the maximum principle alone is insufficient}

The Weitzenb\"ock equation
\[
  0=2\Delta u+4|\nabla_A\Phi|^2+su+u^2,\qquad u=|\Phi|^2,
\]
at a minimum \(x_0\) with \(u(x_0)=0\) gives (using \(\Delta u\ge0\)
at a minimum in the convention \(\Delta=g^{ij}\nabla_i\nabla_j\)):
\[
  2\Delta u(x_0)+4|\nabla_A\Phi|^2(x_0)=0
  \implies|\nabla_A\Phi|^2(x_0)=0,\quad\Delta u(x_0)=0.
\]
So \(\Phi(x_0)=0\) and \(\nabla_A\Phi(x_0)=0\), but this is second-order
vanishing, not infinite-order.  The Dirac equation \(D_A\Phi=0\) does not
force all higher derivatives to vanish: squaring gives
\(\nabla_A^*\nabla_A\Phi+(s/4)\Phi+\text{curvature term}=0\),
and at \(\Phi=0\) only the trace of the Hessian vanishes, not all components.
Standard unique continuation for second-order elliptic operators requires
infinite-order vanishing.

Therefore, isolated zeros of \(|\Phi|\) are possible in principle.  On a
K\"ahler surface, \(\Phi\) corresponds to a holomorphic section with
codimension-2 zero set.  The almost-K\"ahler form \(\omega=\sqrt2\sigma(\Phi)/|\sigma(\Phi)|\) degenerates at these zeros.

\textbf{However}, in the exact equality case (step~6 above), the H\"older
chain in Proposition~3.8 forces \(|\Phi|^4=\text{const}\) directly, giving
\(|\Phi|>0\) everywhere.  The maximum principle is not needed for this step.

\section{Non-collapsing: the product hyperbolic counterexample and the Weyl filter}

\subsection{The counterexample}

On \(M=\Sigma_g\times\Sigma_{g'}\) (products of hyperbolic surfaces),
the degenerating hyperbolic metric \(h_i\) on \(\Sigma_g\) (short geodesic
\(\ell_i\to0\)) gives:
\[
  \operatorname{Ric}_{h_i\oplus k}=-(h_i\oplus k),\qquad
  \operatorname{Vol}(M)=16\pi^2(g-1)(g'-1),
\]
but \(\operatorname{Vol}(B(p_i,r))\le C(r)\ell_i\to0\) near the short
geodesic.  So exact \(\operatorname{Ric}=-g\) with fixed topology and
volume does NOT prevent local collapse.

\subsection{The Weyl filter}

The product \(\Sigma_g\times\Sigma_{g'}\) is K\"ahler with
\(|W_+|^2=s^2/24=2/3\), hence \(|W_+|=2/\sqrt6\).
This is excluded by the Weyl profile: any region with
\(|W_+|\approx2/\sqrt6\) has volume fraction \(O(\delta^2)\to0\).

More generally:
\begin{itemize}
\item Products of surfaces (K\"ahler or anti-K\"ahler):
  \(|W_+|=2/\sqrt6\) or \(0\). Both excluded by profile.
\item Hyperbolic cusps (\(dt^2+e^{-2t/\sqrt3}g_{T^3}\)):
  conformally flat, \(W_+=0\). Excluded by profile (deficit \(1/24\) per
  unit volume).  However, cusps have vanishing volume, so the total
  deficit is small.
\item Nil-geometry collapses: \(W_+\ne2/\sqrt6\) in general, but need
  specific analysis.
\end{itemize}

\textbf{Key observation}: the Weyl profile does NOT directly give
non-collapsing, but it \emph{does} exclude the known models of collapse
under \(\operatorname{Ric}=-g\) as the macroscopic (positive-volume) bulk.
The remaining possibility is collapse in vanishing-volume regions (cusps,
thin necks) that pay negligible profile cost.

\section{Positive branch analysis}

For \(\operatorname{Ric}\approx+g\):
\begin{itemize}
\item Myers: \(\operatorname{diam}(M,g_i)\le\pi\sqrt{3/(1-\delta_i)}\).
\item If \(\operatorname{Vol}(M,g_i)\ge V_*>0\):
  Bishop--Gromov gives \(\operatorname{Vol}(B(p,1))\ge c(D)V_*>0\),
  hence non-collapsing.
\item Anderson compactness \(\to\) smooth Einstein limit.
\end{itemize}

\textbf{Remaining gap for positive branch}: need
\(\operatorname{Vol}(M,g_i)\ge V_*>0\).  For manifolds with monopole
classes (LeBrun type): positive Ricci forces \(s>0\), which gives
\(\Phi\equiv0\) by SW Weitzenb\"ock, contradicting the monopole condition.
So the positive branch is \emph{empty} on these manifolds (no metrics with
tight positive pinching exist at all).

For manifolds without monopole obstructions (e.g., \(\#_k\mathbb{CP}^2\)):
the volume lower bound is an open question.  CGB gives
\(8\pi^2\chi\ge[(2/3)(1-\delta)^2-2\delta^2]\operatorname{Vol}\), which
is an \emph{upper} bound on volume, not a lower bound.

\section{Summary: precise gaps and research priorities}

\textbf{Negative branch on critical Ishida--LeBrun families:}
\begin{enumerate}
\item Weyl profile: rigorously excludes KE bulk, product surface collapse,
  and restricts bubbles to hyperk\"ahler.
\item Armstrong--LeBrun contradiction: rigorous if smooth convergence holds.
\item Gap: non-collapsing (and no-bubble).  All known collapse models are
  excluded by the Weyl profile, but unknown models are not ruled out.
  The specific candidates to exclude are:
  \begin{itemize}
  \item Collapse with \(|W_+|\approx1/(2\sqrt6)\) and
    \(\operatorname{Ric}\approx-g\) in the thin part.
  \item Vanishing-volume cusp/neck regions carrying monopole phase
    transitions and \(W_-\) energy concentration.
  \end{itemize}
\end{enumerate}

\textbf{Positive branch (all dimensions):}
\begin{enumerate}
\item Myers gives diameter bound, Bishop--Gromov gives non-collapsing
  if volume is bounded below.
\item Gap: volume lower bound on fixed manifold with positive Ricci pinching.
\item On LeBrun/SW manifolds: positive branch is \emph{empty} (SW contradiction
  under positive scalar curvature).
\end{enumerate}

\textbf{Most promising next directions:}
\begin{enumerate}
\item Show that the Weyl profile constraint, combined with the
  near-equality SW structure, forces a ``capacity lower bound'' on any
  separating core: that is, the thin regions carrying the topology
  transition must have a minimum volume or Sobolev capacity, preventing
  collapse.
\item Alternatively: study the specific family of metrics with
  \(\operatorname{Ric}\approx-g\) and \(|W_+|\approx1/(2\sqrt6)\) and
  show that no compatible collapse exists (e.g., by analyzing the
  Bianchi identity \(\nabla\cdot W_+=d\operatorname{Ric}^{\circ\#}\)
  and the near-Einstein condition).
\item For the positive branch: establish volume lower bounds using
  topological arguments (Euler characteristic, signature) or collapsing
  theory specific to \(\operatorname{Ric}\ge(1-\delta)g\).
\end{enumerate}
\end{document}
